\input{preamble}

\title{In-Class 3.4 --- Completing the Square --- Answer Key}

\NewDocumentCommand{\ParabVertex}{m m O{2} O{2}}{%
\begin{GridSquareFitEOFilled}[10](2in)
\foreach \x in {-1,-0.5,0,0.5,1} \fill (\x*#4+#1,#3*#4*\x*\x+#2) circle (0.25);
\ifthenelse{#3>0}{
\draw[<->,thick] ({#1-sqrt((#4)/(#3)*(10.9-(#2)))},10.9) parabola bend 
								(#1,#2) ({#1+sqrt((#4)/(#3)*(10.9-(#2)))},10.9);
}{
\draw[<->,thick] ({#1-sqrt((#4)/(#3)*(-10.9-(#2)))},-10.9) parabola bend 
								(#1,#2) ({#1+sqrt((#4)/(#3)*(-10.9-(#2)))},-10.9);
}
\end{GridSquareFitEOFilled}
}
%                                 x x y x x y h k rounded
%                                 1 2 3 4 5 6 7 8 9
\NewDocumentCommand{\ParabPoints}{m m m m m m m m O{#8}}{%
\begin{GridSquareFitEOFilled}[10](2in)
\foreach \x/\y in {#1/#3,#2/#3,#4/#6,#5/#6,#7/#8} \fill (\x,\y) circle (0.25);
\ifthenelse{#9<#6}{ %parabola opens up
\draw[<->,thick] ({#7-sqrt(pow(#1-#7,2)/(#3-#8)*(10.9-#8)},10.9) parabola bend 
								(#7,#8) ({#7+sqrt(pow(#1-#7,2)/(#3-#8)*(10.9-#8)},10.9);
}{
\draw[<->,thick] ({#7-sqrt(pow(#1-#7,2)/(#3-#8)*(-10.9-#8)},-10.9) parabola bend 
								(#7,#8) ({#7+sqrt(pow(#1-#7,2)/(#3-#8)*(-10.9-#8)},-10.9);
}
\end{GridSquareFitEOFilled}
}

\begin{document}
\maketitle

\section*{Completing the Square}
Fill in the blanks to complete the square. Use the squares provided.
\begin{smenum}
\mitemxx
{%
\CS{+8}[16][4]%
\CSsquare[4][16]%
}{%
\CS{-4}[4][(-2)]%
\CSsquare[-2][4]%
}
\mitemxx
{%
\CS{-7}[\frac{49}{4}][\frac{7}{2}]%
\CSsquare[\frac{7}{2}][\frac{49}{4}]%
}{%
\CS{+\frac{7}{5}}[\frac{49}{100}][\frac{7}{10}]%
\CSsquare[\frac{7}{10}][\frac{49}{100}]%
}
\end{smenum}

\section*{Solve by Completing the Square}
Solve each quadratic equation by completing the square.
\begin{smenum}
\mitemxxx
{$\begin{aligned}[t]
x^2-4x+4&=8+4\\
(x-2)^2&=12\\
x-2&=\pm\sqrt{12}\\
x-2&=\pm 2\sqrt 3\\
x&=2\pm 2\sqrt 3
\end{aligned}$}
{$\begin{aligned}[t]
x^2+3x+\frac{9}{4}&=5+\tfrac{9}{4}\\
\left(x+\tfrac{3}{2}\right)^2&=\tfrac{29}{4}\\
x+\tfrac{3}{2}&=\pm\sqrt{\tfrac{29}{4}}\\
x+\tfrac{3}{2}&=\pm\tfrac{\sqrt{29}}{2}\\
x&=-\tfrac{3}{2}\pm\tfrac{29}{2}\\
x&=\tfrac{-3\pm 29}{2}
\end{aligned}$}
{$\begin{aligned}[t]
x^2-4x&=-3\\
x^2-4x+4&=-3+4\\
(x-2)^2&=1\\
x-2&=\pm\sqrt{1}\\
x&=2\pm 1\\
x&=1,\ x=3
\end{aligned}$}
\mitemxxx
{$\begin{aligned}[t]
x^2-7x&=16\\
x^2-7x+\tfrac{49}{4}&=16+\tfrac{49}{4}\\
\left(x-\tfrac{3}{2}\right)^2&=\tfrac{113}{4}\\
x-\tfrac{3}{2}&=\pm\sqrt{\tfrac{113}{4}}\\
x-\tfrac{3}{2}&=\pm\tfrac{\sqrt{113}}{2}\\
x&=\tfrac{3}{2}\pm\tfrac{\sqrt{113}}{2}\\
x&=\tfrac{3\pm\sqrt{113}}{2}
\end{aligned}$}
{$\begin{aligned}[t]
2x^2+8x&=10\\
x^2+4x&=5\\
x^2+4x+4&=5+4\\
(x+2)^2&=9\\
x+2&=\pm\sqrt{9}\\
x+2&=\pm 3\\
x&=-2\pm 3\\
x&=-5,\ x=1
\end{aligned}$}
{$\begin{aligned}[t]
-\tfrac{2}{3}x^2+6x&=-4\\
x^2-9x&=6\\
x^2-9x+\tfrac{81}{4}&=6+\tfrac{81}{4}\\
\left(x-\tfrac{9}{2}\right)^2&=\tfrac{105}{4}\\
x-\tfrac{9}{2}&=\pm\sqrt{\tfrac{105}{4}}\\
x-\tfrac{9}{2}&=\pm\tfrac{\sqrt{105}}{2}\\
x&=\tfrac{9}{2}\pm\tfrac{\sqrt{105}}{2}\\
x&=\tfrac{9\pm\sqrt{105}}{2}
\end{aligned}$}
\end{smenum}

\clearpage
\section*{Graphing Parabolas in Vertex Form}
For each quadratic function below:
\begin{clist}
\item Identify and plot the vertex.
\item Plot two symmetric points with integer coordinates on each side of the vertex.
\end{clist}
\begin{smenum}
\mitemxx
{\ParabVertex{3}{2}}
{\ParabVertex{-1}{-5}[-2][1]}
%{$f(x)=(x-3)^2+2$\\\GridSquareFitEO[10](2in)}
%{$f(x)=-2(x+1)^2-5$\\\GridSquareFitEO[10](2in)}
\mitemxx
{\ParabVertex{2}{-2}[-1][3]}
{\ParabVertex{-5}{-1}[2][4]}
%{$f(x)=-\frac{1}{3}(x-2)^2-2$\\\GridSquareFitEO[10](2in)}
%{$f(x)=\frac{1}{2}(x+5)^2-1$\\\GridSquareFitEO[10](2in)}
\end{smenum}
\clearpage
\section*{Graphing by Rewriting in Vertex Form}
For each quadratic function below:
\begin{clist}
\item Rewrite the function in vertex form.
\item Identify and plot the vertex.
\item Plot two symmetric points with integer coordinates on each side of the vertex.
\end{clist}
\begin{smenum}
\mitemxx
{$\begin{aligned}[t]
f(x)&=x^2-6x\\
&=x^2+6x+9-9\\
&=(x+3)^2-9
\end{aligned}$\\
\ParabVertex{-3}{-9}}
{$\begin{aligned}[t]
f(x)&=-x^2+4x+2\\
&=-\left(x^2-4x\right)+2\\
&=-\left(x^2-4x+4-4\right)+2\\
&=-\left((x-2)^2-4\right)+2\\
&=-(x-2)^2+4+2\\
&=-(x-2)^2+6
\end{aligned}$\\
\ParabVertex{2}{6}[-2][2]
}
\mitemxx
{$\begin{aligned}[t]
f(x)&=2x^2+6x+5\\
&=2\left(x^2+3x\right)+5\\
&=2\left(x^2+3x+\tfrac{9}{4}-\tfrac{9}{4}\right)+5\\
&=2\left(\left(x+\tfrac{3}{2}\right)^2-\tfrac{9}{4}\right)+5\\
&=2\left(x+\tfrac{3}{2}\right)^2-\tfrac{9}{2}+5\\
&=2\left(x+\tfrac{3}{2}\right)^2+\tfrac{1}{2}
\end{aligned}$
\ParabPoints{-3}{0}{5}{-2}{-1}{1}{-1.5}{0.5}[0]
}{$\begin{aligned}[t]
f(x)&=-\frac{1}{2}x^2-3x-2\\
&=-\frac{1}{2}\left(x^2-6x\right)-2\\
&=-\frac{1}{2}\left(x^2-6x+9-9\right)-2\\
&=-\frac{1}{2}\left((x-3)^2-9\right)-2\\
&=-\frac{1}{2}(x-3)^2+\frac{9}{2}-2\\
&=-\frac{1}{2}(x-3)^2+\frac{5}{2}
\end{aligned}$
\ParabPoints{-6}{0}{-2}{-4}{-2}{2}{-3}{2.5}[3]
}
%{$f(x)=-\frac{1}{2}x^2-3x-2$\\Hint: Your points should be $1$ and $3$ units away from the vertex.\\[1ex]\GridSquareFitEO[10](2in)}
\end{smenum}

\end{document}